\section{Open questions} \label{XIV.6}

\subsection{} \label{XIV.6.1} One may ask whether the implication \textup{(ii)} \ALORS \textup{(i)} of \Ref{XIV.4.2} is valid more generally for torsion sheaves $F$, not necessarily annihilated by a given integer $m$ and not necessarily constructible. In the case where $S$ is not of characteristic zero, it seems possible that this implication remains valid, even for sheaves of $p$-torsion ($p$ the residual characteristic). Finally it is also not clear that the hypothesis that $S$ be excellent cannot be removed.

\subsection{} \label{XIV.6.2}
Let $X$ be a scheme proper over a field $k$ or else the complement of the closed point of a henselian local scheme, and $j:Y \to X$ a closed subscheme of $X$, whose complement $U$ is affine. Then, if $F$ is a sheaf of sets on $X$ or a sheaf of not necessarily commutative groups, the statements \Ref{XIV.4.5} or \Ref{XIV.4.6} and \Ref{XIV.4.9} still make sense for such an $F$, provided one restricts oneself to small values of $n$. If $u$ is a point of $U$, one denotes by $\bar{u}$ a geometric point over~$u$, by $X(\bar{u})$ the strict localization of $X$ at $\bar{u}$ and by $F_{\bar{u}}$ the fibre of $F$ at $\bar{u}$. Then, possibly making certain hypotheses on $X$ and on $F$, for example assuming $X$ excellent (possibly of characteristic zero, or of equal characteristic using resolution of singularities) and $F$ ind-finite (or if need be even $\LL$-ind-finite with $\LL$ prime to the characteristic of $X$), one would like to prove the following statements:
\begin{enumeratea}
\item
Let $F$ be a sheaf of sets (\resp a sheaf of groups) and suppose that, for every point $u$ of $U$, one has
$$F_{\bar{u}} \to H^0(X(\bar{u})-\bar{u}, F)\text{ injective if } \dim(\overline{\{u\}}) \leq 1$$
(that is, for such a $u$, one has $\prof_u(F) \geq 1$). Then, as $V$ runs through the set of open neighbourhoods of $Y$, the canonical morphism
$$
\varinjlim_V H^0(V, F) \to H^0(Y, j^*F)
$$
is bijective (\resp one has the preceding conclusion and moreover the morphism $\varinjlim_V H^1(V, F) \to H^1(Y, j^* F)$ is injective). If $F$ is constructible, one may replace the $\varinjlim$ by the cohomology of $V$ for $V$ \og small enough\fg.
\item
Let $F$ be a sheaf of sets (\resp a sheaf of groups) and suppose that, for every point $u$ of $U$, one has $\prof_u(F) \geq 2 -\dim(\overline{\{u\}})$, which is also expressed by the relations
\begin{gather*}
F_{\bar{u}}\to H^0(X(\bar{u})-\bar{u}, F)\text{ is bijective if } \dim(\overline{\{u\}})=0\\
F_{\bar{u}}\to H^0(X(\bar{u})-\bar{u}, F)\text{ is injective if } \dim(\overline{\{u\}})=1.
\end{gather*}
Then the canonical morphism
$$H^0(X, F) \to H^0(Y, j^*F)$$
is bijective (\resp one has the preceding conclusion and moreover the morphism $H^1(X, F) \to H^1(Y, j^*F)$ is injective).
\item
Let $F$ be an ind-finite sheaf of groups. Suppose that, for every point $u$ of $U$, one has
\begin{gather*}
F_{\bar{u}}\to H^0(X(\bar{u})-\bar{u}, F)\text{ bijective if } \dim(\overline{\{u\}})=0\text{ or }1,\\
F_{\bar{u}}\to H^0(X(\bar{u})-\bar{u}, F)\text{ injective if } \dim(\overline{\{u\}})=2.
\end{gather*}
Then, as $V$ runs through the set of open neighbourhoods of $Y$, the canonical morphisms
$$
\varinjlim_V H^0(V, F) \to H^0(Y, j^*F) \text{ and } \varinjlim_V H^1(V, F) \to H^1(Y, j^*F)
$$
are bijective. If $F$ is constructible, one may replace the $\varinjlim$ by the cohomology of~$V$ for $V$ small enough.
\item
Let $F$ be a sheaf of groups. Suppose that, for every point $u$ of $U$, one has $\prof_u(F) \geq 3-\dim(\overline{\{u\}})$, which is also expressed by the conditions
\begin{gather*}
F_{\bar{u}}\to H^0(X(\bar{u})-\bar{u}, F)\text{ bijective, and } H^1(X(\bar{u})-\bar{u}, F)=0 \text{ if } \dim(\overline{\{u\}})=0,\\
F_{\bar{u}}\to H^0(X(\bar{u})-\bar{u}, F)\text{ bijective if } \dim(\overline{\{u\}})=1,\\
F_{\bar{u}}\to H^0(X(\bar{u})-\bar{u}, F)\text{ injective if } \dim(\overline{\{u\}})=2.
\end{gather*}
Then the canonical morphisms
$$H^0(X, F) \to H^0(Y, j^*F)\text{ and } H^1(X, F) \to H^1(Y, j^*F)$$
are bijective.
\end{enumeratea}

As an indication in favour of these statements\nde{\label{MiR}all the statements of~\Ref{XIV.6.2}, apart from the constructible variants, have been proved by Mme Raynaud; see (Raynaud~M., {\og Th\'eor\`emes de Lefschetz en cohomologie des faisceaux coh\'erents et en cohomologie \'etale. Application au groupe fondamental\fg}, \emph{Ann. Sci. \'Ec. Norm. Sup. (4)} \textbf{7} (1974), p\ptbl 29--52, corollary III.1.3).}, we mention \Ref{XIII}~\Ref{XIII.2.1}, \Ref{X}~\Ref{X.3.4} and \Ref{XII}~\Ref{XII.3.5}. We mention that, thanks to the argument of \Ref{XIV.4.8} and \Ref{XIV.4.9}, statement a) (\resp c)) would follow from b) (\resp d)).

\subsection{} \label{XIV.6.3}
There would follow from d) the following statement, analogous to \Ref{XIV.5.6}: if $A$ is a noetherian local ring (possibly excellent) and if one has $\prof \geom (A) \geq 3$, then one has $\prof \hop (A) \geq 3$. To see this, one realises $Y'=\Spec A$ as a closed subscheme of a regular local scheme $X'=\Spec B$, whose complement is a union of $q$ affine open sets, with the relation $\dim B -q=\prof \geom(A)$. One has, for every point $x$ of $X'$, if $n=\dim B$, $\prof \hop_x(X) \geq \inf(3, n-\dim(\overline{\{x\}}))$ (\cf \Ref{XIV.1.11}) and one deduces from d) that this implies $\prof \hop_y(Y') \geq \inf(3, n-q-\dim(\overline{\{y\}}))$, for every point $y$ of $Y'$. The result is then obtained by taking for $y$ the closed point of $Y'$.

\subsection{} \label{XIV.6.4}
A variant of \Ref{XIV.4.2}, at least of the implication \textup{(ii)} \ALORS \textup{(i)}, should still be valid in the complex analytic case, provided one works with \og analytically constructible\fg sheaves (\cf \Ref{XIII}); the proof would be analogous to that of \Ref{XIV.4.2}, using the duality theory of J.-L.~Verdier. Note on the other hand that, for the complex analytic analogue of the non-commutative variants indicated in \Ref{XIV.6.2}, one does not even have a method of attack for the statements concerning the fundamental group suggested by the results of Expos\'es \Ref{X}, \Ref{XII}, \Ref{XIII}, recalled at the end of \Ref{XIV.6.2}. The methods of the Seminar seem in fact to be irremediably tied to the case of \emph{finite} coverings (which can be studied in terms of coherent sheaves of algebras).

\begin{thebibliography}{9}
\bibitem[1]{XIV.1}%
\textsc{M\ptbl Artin \& B\ptbl Mazur} -- {\og Homotopy of Varieties in Etale Topology\fg}, in \emph{Proceedings of a Conference on Local Fields}, Springer, 1967.

\bibitem[2]{XIV.2}%
\textsc{J.-P\ptbl Serre} -- \emph{Cohomologie Galoisienne}, Springer-Verlag, 1964.

\bibitem[3]{XIV.3}%
\textsc{J.-L\ptbl Verdier} -- \emph{Des cat\'egories d\'eriv\'ees des cat\'egories ab\'eliennes}, with a preface by Luc Illusie, edited by Georges Maltsiniotis, Ast\'erisque, vol.~239, Soci\'et\'e math\'ematique de France, Paris, 1996.

\end{thebibliography}
